% WWW table style: three rules, bold headers, reference green highlight.
\section{Related Work}
\begin{table*}[t]
\centering
\caption{\textbf{Mechanisms in memory-efficient adaptive optimization.} \yes: present; \no: absent; --: no coordinate-alignment operation. ``No refresh SVD'' refers to full-gradient decompositions after initialization; small tangent decompositions are allowed. Moment alignment uses projection-aware approximations, not exact covariance transport. Generalized error feedback (EF) differs from an immediate residual update.}
\label{tab:related-mechanisms}
\setlength{\tabcolsep}{3.2pt}
\renewcommand{\arraystretch}{1.24}
\small
\reviewfit{\begin{tabular}{@{}l l l c c@{}}
\toprule {\bfseries Method} & {\bfseries Subspace rule} & {\bfseries Outside subspace} & {\bfseries Moment} & {\bfseries No refresh}\\{\bfseries } & {\bfseries } & {\bfseries } & {\bfseries alignment} & {\bfseries SVD}\\
\midrule
GaLore~\cite{zhao2024galore}  &  Periodic SVD  &  Discarded  &  \no  &  \no\\
Fira~\cite{chen2024fira}  &  Periodic SVD  &  Scaled residual \yes  &  \no  &  \no\\
APOLLO~\cite{zhu2024apollo}  &  Random / SVD  &  Full-gradient scaling  &  --  &  \yes$^{a}$\\
LDAdam~\cite{robert2024ldadam}  &  Block power + QR  &  Generalized EF  &  \yes  &  \yes\\
SubTrack++~\cite{rajabi2025subtrackpp}  &  Grassmann geodesic  &  Scaled residual \yes  &  \yes  &  \yes\\
\rowcolor{oursbg}
\textbf{GrassRoute}  &  \textbf{Grassmann QR}  &  \textbf{Norm-matched residual \yes}  &  \yes  &  \yes\\
\bottomrule
\end{tabular}}
\par\vspace{3pt}
\begin{minipage}{\linewidth}\footnotesize
$^{a}$Random-projection APOLLO; its SVD variant uses refresh decompositions.
GrassRoute's analysis covers normalized-covariance geometry, tracking, adaptive convergence, full-state stability, and PAC-Bayes generalization, each under its stated assumptions.
\end{minipage}
\end{table*}

Low-rank optimizers reduce memory by compressing adaptive state. GaLore periodically recomputes a gradient subspace by SVD \cite{zhao2024galore}. FLoRA, Fira, APOLLO, OSD, LDAdam, and related methods vary the projection rule, residual update, or treatment of adaptive statistics \cite{hao2024flora,chen2024fira,zhu2024apollo,liang2024osd,robert2024ldadam,he2024golore,su2025galore2,mo2025loro}. Table~\ref{tab:related-mechanisms} compares the mechanisms most relevant to GrassRoute.

Grassmann geometry treats bases related by right orthogonal rotations as the same subspace \cite{edelman1998geometry,absil2008optimization,boumal2014manopt}. GROUSE and related streaming algorithms use this geometry for online subspace estimation \cite{balzano2011grouse,balzano2018streaming}. The closest optimizer is SubTrack++, which combines Grassmann geodesic updates, projection-aware Adam, and recovery scaling \cite{rajabi2025subtrackpp}. GrassRoute uses a right-sided QR tracker and a norm-matched residual update. Online manifold tracking and moment alignment are established ingredients; our focus is their approximation error, convergence conditions, and empirical memory trade-offs beyond language-model training.

Our proofs build on Rayleigh-quotient geometry \cite{alimisis2024geodesic}, spectral perturbation \cite{yu2015useful}, adaptive-direction analysis \cite{reddi2018adam,chen2019adamtype}, and algorithmic stability \cite{bousquet2002stability,hardt2016stability}. We apply these tools to the normalized population covariance and the full optimizer state, including the residual branch. The generalization analysis uses full-state coupling and, separately, the Seeger--Maurer binary-KL PAC-Bayes bound \cite{seeger2002pacbayes,maurer2004note}. Davis--Kahan perturbation theory and variance-sensitive matrix concentration connect subspace error to posterior complexity \cite{davis1970rotation,tropp2012user}. The proofs state the eigengap, angle, and independence conditions needed for these connections.
